\chapter*{Заключение}                       % Заголовок
\addcontentsline{toc}{chapter}{Заключение}  % Добавляем его в оглавление

%% Согласно ГОСТ Р 7.0.11-2011:
%% 5.3.3 В заключении диссертации излагают итоги выполненного исследования, рекомендации, перспективы дальнейшей разработки темы.
%% 9.2.3 В заключении автореферата диссертации излагают итоги данного исследования, рекомендации и перспективы дальнейшей разработки темы.

В диссертационной работе решена научная задача разработки теоретического подхода к анализу связи между объемом обучающей выборки, локальными свойствами оптимизационной поверхности и достаточным размером данных в параметрических моделях машинного обучения, а также предложена интерпретация полученных результатов в терминах законов масштабирования.

Основные результаты работы заключаются в следующем.
\begin{enumerate}[beginpenalty=10000]
    \item Для нейронных сетей исследованы спектральные свойства матрицы Гессе функции потерь. Показано, что в окрестности точки минимума спектрально значимая часть второй производной естественным образом описывается гаусс"--~ньютоновской компонентой. В классе матрично-представимых нейронных сетей получена ее блочная факторизация, на основе которой выведены архитектурно-зависимые верхние оценки спектральной нормы для полносвязных и сверточных архитектур.

    \item Разработан подход к анализу стабилизации ландшафта функции потерь при увеличении объема обучающей выборки. На основе локальной квадратичной аппроксимации функции потерь в окрестности точки минимума введен общий критерий $p$-сходимости, а также рассмотрены его частные случаи: абсолютный одноточечный критерий, среднеквадратичный критерий в полном пространстве параметров и среднеквадратичный критерий в подпространстве наибольшей кривизны.

    \item Для введенных критериев стабилизации получены теоретические оценки скорости сходимости. Показано, что абсолютный одноточечный критерий убывает с порядком $\mathcal{O}(k^{-1})$, тогда как среднеквадратичный критерий в полном пространстве параметров и его подпространственный аналог убывают с порядком $\mathcal{O}(k^{-2})$ при увеличении объема выборки.

    \item Построены критерии достаточного размера выборки для параметрических моделей машинного обучения, основанные на стабилизации ландшафта функции потерь. Показано, что верхние оценки локальной кривизны оптимизационной поверхности, выраженные через спектральную норму гаусс"--~ньютоновской компоненты, позволяют получить архитектурно-зависимые оценки достаточного объема данных для полносвязных и сверточных моделей.

    \item Предложена теоретическая схема интерпретации связи между архитектурой модели, локальной геометрией оптимизационной поверхности и объемом обучающей выборки в терминах законов масштабирования. Показано, что вклад конечности обучающей выборки в итоговое значение функции потерь может быть связан с локальной кривизной поверхности, а оптимальное соотношение между размером модели и размером выборки зависит не только от ошибки выразимости, но и от режима роста архитектурно-зависимой кривизны.

    \item Для вероятностных линейных моделей исследована стабилизация распределения оптимальных параметров при увеличении объема обучающей выборки. Введены и проанализированы критерии достаточного размера выборки, основанные на дисперсии и математическом ожидании функции правдоподобия, дивергенции Кульбака"--~Лейблера между апостериорными распределениями и метрике близости s-score. Для линейной регрессии с гауссовским априорным распределением получены условия корректности этих критериев и исследовано асимптотическое поведение соответствующих вероятностных характеристик.

    \item Проведенные вычислительные эксперименты для нейросетевых и линейных моделей качественно согласуются с полученными теоретическими результатами. Для нейронных сетей подтверждена близость спектральных норм полной матрицы Гессе и ее гаусс"--~ньютоновской компоненты на поздних этапах обучения, а также наблюдаемая стабилизация введенных критериев при увеличении объема выборки. Для линейных моделей подтверждена сходимость предложенных вероятностных критериев достаточности.
\end{enumerate}

Полученные результаты имеют теоретическое значение для развития математических основ машинного обучения и анализа параметрических моделей, поскольку связывают локальную геометрию оптимизационной поверхности, спектральные свойства матрицы Гессе и достаточный размер обучающей выборки в рамках единой постановки. Практическая значимость работы состоит в возможности использования разработанных критериев и оценок при теоретически обоснованной оценке достаточного объема данных, при анализе чувствительности моделей к увеличению выборки, а также при интерпретации достаточности выборки в терминах стабилизации ландшафта функции потерь.

К числу перспектив дальнейшего развития темы относятся:
\begin{enumerate}[beginpenalty=10000]
    \item распространение полученных результатов на более широкий класс архитектур, включая трансформерные и гибридные нейросетевые модели;
    \item уточнение связи между стабилизацией ландшафта функции потерь и обобщающей способностью моделей;
    \item развитие более точных архитектурно-зависимых оценок локальной кривизны с учетом нормализаций, остаточных связей и современных механизмов внимания;
    \item исследование критериев достаточного размера выборки для обобщенно-линейных и нелинейных вероятностных моделей;
    \item разработка вычислительно эффективных процедур практической оценки достаточности выборки для больших моделей машинного обучения.
\end{enumerate}

Таким образом, в работе предложен теоретический подход к исследованию достаточного размера обучающей выборки на основе анализа локальных свойств оптимизационной поверхности, получены строгие результаты для нейросетевых и линейных моделей и намечены направления дальнейшего развития этой тематики в рамках специальности 1.2.1~--- <<Искусственный интеллект и машинное обучение>>.

В заключение автор выражает благодарность и глубокую признательность научному руководителю к.ф.-м.н.~Грабовому~А.\,В. за поддержку, помощь, обсуждение результатов и научное руководство.